% Problem-section file following ../_template.tex.
\section{Secret key from every entangled state}
\label{sec:problem-20}

\paragraph{Problem.}
Does every finite-dimensional entangled bipartite state have positive
asymptotic distillable secret key?  For a state $\rho_{AB}$, with an adversary
holding a purification, define its distillable-key rate by
\begin{equation}
  K_D(\rho_{AB})
  :=\sup\left\{R\geq0:
    \begin{array}{l}
      \text{there are LOPC protocols $\Lambda_n$ and private states
      $\gamma^{(n)}_{K_n}$, with $K_n=2^{\lfloor nR\rfloor}$, such that}\\[-1mm]
      \displaystyle
      \lim_{n\to\infty}
      \left\|\Lambda_n(\rho_{AB}^{\otimes n})
      -\gamma^{(n)}_{K_n}\right\|_1=0
    \end{array}
  \right\},
  \label{eq:p20-distillable-key}
\end{equation}
where each target $\gamma^{(n)}_{K_n}$ may have its own shield system and has
key registers of dimension $K_n$.  With the operational convention in
Eq.~\eqref{eq:p20-distillable-key},
the question is whether the implication
\begin{equation}
  \rho_{AB}\ \text{entangled}
  \quad\Longrightarrow\quad
  K_D(\rho_{AB})>0
  \label{eq:p20-entanglement-implies-key}
\end{equation}
holds for every finite-dimensional $\rho_{AB}$.

\subsection*{Status}

Unsolved

\subsection*{Source}

Horodecki, Sikorski, Das, and Wilde explicitly identify the question whether
every entangled state has positive distillable key
\sourcecite{ref:p20-key-cost}{HSDW26}.

\subsection*{Progress}

\begin{itemize}
  \item Horodecki, Horodecki, Horodecki, and Oppenheim constructed
  bound-entangled states with positive distillable key.  Thus zero
  distillable entanglement does not provide a counterexample to
  Eq.~\eqref{eq:p20-entanglement-implies-key}
  \sourcecite{ref:p20-bound-key}{HHHO05}.

  \item Horodecki, Sikorski, Das, and Wilde explicitly state that the
  existence of an entangled state with zero distillable key remains open.
  Their 2026 result therefore confirms the unresolved scope of
  Eq.~\eqref{eq:p20-entanglement-implies-key} close to the present cutoff
  \sourcecite{ref:p20-key-cost}{HSDW26}.
\end{itemize}

\subsection*{References}

\begin{enumerate}[label={},leftmargin=4.8em,labelsep=0.5em]
  \footnotesize
  \item[\textup{[HHHO05]}]\label{ref:p20-bound-key}
  K. Horodecki, M. Horodecki, P. Horodecki, and J. Oppenheim,
  ``Secure Key from Bound Entanglement,''
  \emph{Physical Review Letters} \textbf{94}, 160502 (2005).
  \href{https://doi.org/10.1103/PhysRevLett.94.160502}{doi:10.1103/PhysRevLett.94.160502};
  \href{https://arxiv.org/abs/quant-ph/0309110}{arXiv:quant-ph/0309110}.

  \item[\textup{[HSDW26]}]\label{ref:p20-key-cost}
  K. Horodecki, L. Sikorski, S. Das, and M. M. Wilde,
  ``Cost of Quantum Secret Key,'' \emph{Quantum} \textbf{10}, 2098 (2026).
  \href{https://doi.org/10.22331/q-2026-05-06-2098}{doi:10.22331/q-2026-05-06-2098};
  \href{https://arxiv.org/abs/2402.17007}{arXiv:2402.17007}.
\end{enumerate}

\subsection*{Comment}

The unresolved alternative is either to prove
Eq.~\eqref{eq:p20-entanglement-implies-key} or to construct an entangled state
with $K_D=0$.  The question concerns the asymptotic key rate, not one-shot key
extraction or distillable entanglement.

\subsection*{Tag}

Quantum cryptography; Bound entanglement; Entanglement theory;
Local operations and classical communication

\subsection*{ID}

\texttt{op\_b952ec2dd8246a10}
